\documentclass[12pt]{article}
\usepackage[utf8]{inputenc}
\usepackage[T1]{fontenc}
\usepackage[margin=1in]{geometry}
\usepackage{setspace}
\usepackage{xcolor}
\usepackage[colorlinks=true, urlcolor=blue!60!black, linkcolor=blue!60!black]{hyperref}
\setlength{\parindent}{0em}
\setlength{\parskip}{0.6em}
\renewcommand{\baselinestretch}{1.0}
\begin{document}
\begin{center}
{\LARGE \textit{Taylor v. United States}, 579 U.S. 301 (2016)\par}
\vspace{0.5em}
{\large Case Brief --- Two-Page Edition\par}
\end{center}

\section{Memory Jogger}
Robbing or attempting to rob a drug dealer of drugs or drug proceeds satisfies the Hobbs Act's commerce element as a matter of law under \textit{Gonzales v. Raich} (\href{https://www.law.cornell.edu/supremecourt/text/14-6166}{Cornell LII}).

\section{Detailed Case Facts}
David Taylor, a member of Roanoke's ``Southwest Goonz,'' joined two 2009 home invasions targeting marijuana dealers, assaulting residents and leaving with jewelry, \$40, and cell phones. He was charged with Hobbs Act robbery and a firearm offense.

\section{Procedural History}
After a hung jury, the district court barred evidence that the dealers sold only local marijuana; Taylor was convicted and sentenced to 28 years. The Fourth Circuit affirmed, splitting with the Second and Seventh Circuits.

\section{Judicial Votes}
Affirmed, 7--1. Alito for the Court (Roberts, Kennedy, Ginsburg, Breyer, Sotomayor, Kagan); Thomas dissenting.

\section{Holding}
The government need not prove the drugs crossed state lines. Proof that the defendant knowingly stole or tried to steal drugs or drug proceeds suffices, because under \textit{Raich} the whole drug market is ``commerce over which the United States has jurisdiction.'' The holding is limited to robberies targeting drug dealers.

\section{Analysis of Opinions}
\textbf{Alito:} The Hobbs Act is ``unmistakably broad,'' and \textit{Raich} held that drug production and sale, even local, is a regulable class of activity. The jury finds facts beyond a reasonable doubt, but the meaning of the commerce element is a question of law. \textbf{Thomas:} Robbery is not commerce; the substantial-effects test has ``no stopping point,'' expands \textit{Raich}, and relieves prosecutors of proving a real nexus.

\section{Examples: Future Applications}
\textbf{Same-Side:} Robbing a local methamphetamine dealer of proceeds satisfies the element. \textbf{Opposite-Side:} Burglarizing a retired teacher's home with no drug connection requires proof of an actual nexus. \textbf{Fence-Sitter:} Robbing a state-licensed intrastate cannabis dispensary raises an open question.

\section{Critique}
Critics echo Thomas that \textit{Taylor} federalizes local crime, and reformers decry a 28-year sentence for a ``bit of a tagalong.'' Scholars question whether a deference precedent can establish an element of a crime.

\section{Key Quotations}
\begin{quote}
``Rather, to satisfy the Act's commerce element, it is enough that a defendant knowingly stole or attempted to steal drugs or drug proceeds, for, as a matter of law, the market for illegal drugs is `commerce over which the United States has jurisdiction.' ''
\end{quote}
\noindent --- Justice Alito, opinion of the Court, 579 U.S. at 309.

\end{document}
